\section{Discussion and Conclusion}
\label{sec:discussion}

The proposed Digitally Assisted Analog Front-End architecture demonstrates that embedded mixed-signal resources can improve the effective operating conditions of configurable analog hardware without modifying the analog acquisition path during normal operation. Unlike digital post-processing, the embedded resources temporarily assist the analog circuitry during foreground calibration and are subsequently disengaged, allowing the calibrated signal chain to operate as a conventional feed-forward analog front end.

From an instrumentation perspective, the proposed architecture differs from two conventional approaches. Unlike precision analog front ends designed to operate close to the sensor's Johnson-noise limit \cite{Kirchhoff2017}, it trades analog component precision for embedded measurement and digital assistance. Likewise, unlike purely digital offset compensation, it physically recenters observable additive DC errors within the available VDAC authority and range instead of correcting them numerically after acquisition. The reported improvements therefore reflect the performance of the complete architecture using a common fixed-reference baseline and a single controller configuration across all evaluated board and gain combinations rather than per-stage parameter optimization.

The proposed foreground calibration strategy also differs from continuous DC-servo architectures, which inherently introduce a settling-versus-bandwidth tradeoff through permanently closed feedback loops \cite{Liu2020}. Here, digital assistance is active only during foreground calibration. Once calibration is completed, the controller is disabled, the VDAC outputs remain constant, and the acquisition path contains no continuously active digital feedback.

\label{sec:pga-note}
The present implementation should be interpreted as a proof of architectural feasibility rather than as an optimized analog front end. The relatively modest PGA improvement illustrates this point. Because $|K_{\mathrm{PGA}}|=1$, the selected deadband ($\Delta_{\mathrm{PGA}}=18$) accepts integer-error buckets through approximately $\pm288$~mV (continuous edge near $\pm296$~mV after rounding). Consequently, the 0.24-V residual shown in Fig.~\ref{fig:calibration_trace} reflects the selected stopping criterion rather than controller instability. A separate diagnostic experiment on Board~B at $\times4$ gain, performed outside the reported comparison, reduced the LP residual from 2.98 to 0.84\%FS using a narrower deadband, indicating that further improvement is possible through parameter optimization. Likewise, the only increase observed at the reported precision (Board~A, SUM stage, $\times50$ gain, from 0.4 to 0.8\%FS) results from deadband acceptance and does not propagate to the downstream LP output.

The timing results also support the practical feasibility of foreground calibration. Visible correction in Fig.~\ref{fig:calibration_trace} spans approximately 7~s, while informal bench observations indicated completion in about 10~s. The configured timeout windows total 57.6~s (approximately 60.5~s including settling and refinement), excluding stored-code verification and switching overhead. No timeout occurred during the reported experiments; the only timeout observed in bench testing was a board powered while its geophone was moving, whose AC transient violated the quasi-static assumption required during foreground calibration. Nevertheless, the present study was intentionally limited to two independently assembled boards and fixed calibration ordering; consequently, it does not establish population-level statistical performance, long-term drift behavior, or optimal controller tuning.

Overall, the proposed Digitally Assisted Analog Front-End architecture offers a practical means of improving the operating conditions of low-cost configurable analog front ends by temporarily exploiting embedded mixed-signal resources. The reported foreground DC-offset calibration consistently reduced the residual offset across all evaluated boards and gain settings using a single controller configuration, while preserving a conventional analog acquisition path. Although demonstrated through DC-offset calibration, the architecture is not restricted to this function and can be extended to other digitally assisted tasks in future work.
